\begin{ouverture}[avancement]{8}{Réactions chimiques et bilans}{Constitution et transformations}{Observer · Modéliser · Expérimenter · Expliquer}
\objectif{Identifier réactifs, produits et espèces spectatrices.}
\objectif{Ajuster une équation en conservant atomes et charge.}
\objectif{Déterminer le réactif limitant et interpréter un effet thermique.}
\prerequis{Proportionnalité, conversions et lecture de graphiques.}
\end{ouverture}
\essentiel
\essSection{1}{Observer puis modéliser}
Une transformation chimique modifie les espèces présentes : certains réactifs sont consommés et des produits apparaissent. Un dégagement gazeux ou une couleur qui change est un indice à compléter par l'identification des espèces. L'équation ne décrit pas toutes les étapes microscopiques ; elle fournit un bilan.
\begin{savoirfaire}[Ajuster]Écrire les formules correctes, puis changer uniquement les coefficients devant les formules. Conserver le nombre d'atomes de chaque élément et la charge totale. Pour une combustion complète du méthane : $\ce{CH4 + 2 O2 -> CO2 + 2 H2O}$.\end{savoirfaire}
Les espèces spectatrices sont présentes mais inchangées ; on peut les omettre de l'équation ionique simplifiée.
\essSection{2}{Proportions et réactif limitant}
Les coefficients donnent des proportions de nombres d'entités ou de quantités de matière, jamais directement de masses. Dans $\ce{2 H2 + O2 -> 2 H2O}$, deux moles de dihydrogène réagissent avec une mole de dioxygène.
\schema[corrige]{reaction}
\begin{definition}[Transformation totale]Le réactif limitant est entièrement consommé à l'issue de la transformation totale modélisée. Pour le trouver, comparer les quantités disponibles divisées par leurs coefficients ; le plus petit rapport fixe la limite. Si les rapports sont égaux, les proportions sont stœchiométriques.\end{definition}
La présence d'un réactif dans l'état final peut montrer qu'il était en excès, si l'on sait que la transformation a été totale.
\essSection{3}{Réactions usuelles et effet thermique}
Exemples à savoir exploiter : $\ce{C + O2 -> CO2}$ ; $\ce{Zn + 2 H+ -> Zn^{2+} + H2}$ ; $\ce{CaCO3 + 2 H+ -> Ca^{2+} + CO2 + H2O}$ ; $\ce{H+ + HO- -> H2O}$. Les ions chlorure d'une solution d'acide chlorhydrique peuvent rester spectateurs.
\begin{propriete}[Échanges d'énergie]Une transformation exothermique cède de l'énergie au milieu ; une transformation endothermique en reçoit. Une hausse de température dans un montage suffisamment isolé est un indice du caractère exothermique. La température dépend aussi des pertes et des quantités utilisées.\end{propriete}
\essSection{4}{Expérimenter et modéliser}
\begin{experience}[Repérer le réactif limitant]Au laboratoire encadré, faire réagir de petites quantités de carbonate de calcium avec une solution acide ; mettre en évidence le gaz formé et rechercher du solide restant. Comparer des essais en faisant varier une seule quantité initiale. Ne pas fermer hermétiquement un montage qui dégage un gaz.\end{experience}
\begin{experience}[Suivre la température]Relever la température avant mélange puis son évolution, avec le même récipient et le même volume total pour comparer plusieurs essais. Doubler la masse du réactif limitant n'implique une proportionnalité de l'effet observé que si les autres conditions restent adaptées.\end{experience}
\exercices
\rubrique{Prendre confiance}
\begin{blocLivre}{ex-8-1}\begin{exercice}[1]{Lire une équation}
$\ce{C + O2 -> CO2}$.
\begin{questions}
\item Identifier réactifs et produit.
\item Vérifier la conservation des atomes.
\end{questions}
\end{exercice}\end{blocLivre}
\begin{blocLivre}{ex-8-2}\begin{exercice}[1]{Ajuster l’eau}
On propose $\ce{H2 + O2 -> H2O}$.
\begin{questions}
\item Ajuster les coefficients.
\item Pourquoi ne faut-il pas remplacer \ce{H2O} par \ce{H2O2} ?
\end{questions}
\end{exercice}\end{blocLivre}
\begin{blocLivre}{ex-8-3}\begin{exercice}[1]{Un métal et un acide}
Le zinc réagit selon $\ce{Zn + 2 H+ -> Zn^{2+} + H2}$.
\begin{questions}
\item Quelle espèce gazeuse est produite ?
\item Vérifier la charge totale.
\end{questions}
\end{exercice}\end{blocLivre}
\begin{blocLivre}{ex-8-4}\begin{exercice}[1]{Un effet thermique}
La température d’un mélange passe de 20,0 à 26,5 °C dans un récipient approximativement isolé.
\begin{questions}
\item Quel caractère thermique est suggéré ?
\item Quelle grandeur expérimentale a-t-on suivie ?
\end{questions}
\end{exercice}\end{blocLivre}
\rubrique{S’entraîner}
\begin{blocLivre}{ex-8-5}\begin{exercice}[2]{Un premier limitant}
Pour $\ce{2 H2 + O2 -> 2 H2O}$, on dispose de 0,60 mol de \ce{H2} et 0,20 mol de \ce{O2}.
\schema{reaction}
\begin{questions}
\item Déterminer le réactif limitant.
\item Calculer la quantité d’eau formée et le dihydrogène restant.
\end{questions}
\end{exercice}\end{blocLivre}
\begin{blocLivre}{ex-8-6}\begin{exercice}[2]{Espèces spectatrices}
On mélange une solution contenant \ce{H+}, \ce{Cl-} et une solution contenant \ce{Na+}, \ce{HO-}. Les ions \ce{Na+} et \ce{Cl-} restent inchangés.
\begin{questions}
\item Nommer les espèces spectatrices.
\item Écrire l’équation simplifiée.
\end{questions}
\end{exercice}\end{blocLivre}
\begin{blocLivre}{ex-8-7}\begin{exercice}[2]{Le méthane brûle}
Une combustion complète consomme 0,150 mol de méthane avec suffisamment de dioxygène.
\begin{questions}
\item Écrire l’équation ajustée.
\item Donner les quantités de dioxygène consommé et de dioxyde de carbone produit.
\end{questions}
\end{exercice}\end{blocLivre}
\begin{blocLivre}{ex-8-8}\begin{exercice}[2]{Un solide reste}
Après réaction totale entre du calcaire et un acide, du calcaire reste au fond du récipient.
\begin{questions}
\item Quel réactif est en excès ?
\item Quel réactif a limité la transformation, dans le modèle à deux réactifs ?
\end{questions}
\end{exercice}\end{blocLivre}
\clearpage
\rubrique{Raisonner et approfondir}
\begin{blocLivre}{ex-8-9}\begin{exercice}[3]{Détartrer avec des proportions}
$\ce{CaCO3 + 2 H+ -> Ca^{2+} + CO2 + H2O}$. Quantités initiales : 0,080 mol de calcaire et 0,100 mol de \ce{H+}.
\begin{questions}
\item Déterminer le limitant.
\item Calculer le calcaire restant et le \ce{CO2} formé.
\end{questions}
\end{exercice}\end{blocLivre}
\begin{blocLivre}{ex-8-10}\begin{exercice}[3]{Conserver une masse totale}
Dans un système fermé, 4,0 g de dihydrogène réagissent totalement avec 32,0 g de dioxygène et ne produisent que de l’eau.
\begin{questions}
\item Calculer la masse d’eau produite.
\item Pourquoi une expérience ouverte pourrait-elle sembler perdre de la masse ?
\end{questions}
\end{exercice}\end{blocLivre}
\begin{blocLivre}{ex-8-11}\begin{exercice}[3]{Une proportion à ne pas confondre}
Pour $\ce{2 H2 + O2 -> 2 H2O}$, un élève affirme qu’il faut deux grammes de \ce{H2} pour un gramme de \ce{O2}.
\begin{questions}
\item Pourquoi son interprétation est-elle fausse ?
\item Donner la proportion correcte avec les quantités de matière.
\end{questions}
\end{exercice}\end{blocLivre}
\begin{blocLivre}{ex-8-12}\begin{exercice}[3]{Comparer des essais thermiques}
Deux essais utilisent le même récipient, le même volume total et le même excès d’un réactif. La quantité du réactif limitant est doublée.
\begin{questions}
\item Quel effet peut-on attendre sur l’énergie libérée si la réaction totale reste la même ?
\item Pourquoi la hausse de température mesurée peut-elle s’écarter du double ?
\end{questions}
\end{exercice}\end{blocLivre}

